\clearpage
\section{Common allowances and observation transfer}
\label{sec:copying-observation-padding}
The preceding section proves structural filling for all original term contexts.
To use that correspondence semantically, comparisons made at different finite
representation allowances must be related. This section proves that every
closing context lifts through padding, that observational refinement at a
larger allowance reflects to the original allowance, and that the existing
all-allowances definition of term equivalence is equivalent to checking a
cofinal collection of common allowances.

The structural statements are unconditional. The semantic transfer theorems
explicitly require execution invariance under recursive graph isomorphism and
one-step padding of closed programs. They do not assume determinism or template
robustness. These representation laws remain obligations for the copying UAM;
they are not project axioms. A checked counterexample below shows that even
determinism and isomorphism invariance do not imply padding invariance.

\subsection{Exact padding of programs and closing contexts}

Padding a graph at allowance $d$ to $d+n$ equals its original $n$-step
padding. The proof transports the equality $(d+n)-d=n$ through the dependent
graph index; the resulting equality of depth proofs is immaterial.

Choosing an intermediate sufficient allowance changes nothing:
\[
 \mathsf{padTo}_{N}(\mathsf{padTo}_{e}(G))
   =\mathsf{padTo}_{N}(G)\qquad(d\leq e\leq N).
\]
Write $e=d+n$ and $N=(d+n)+m$. Reduce the two successive paddings to their
iterated forms, apply the previous section's additive padding equality, and
transport the final allowance along associativity of addition. Equality
elimination shows that transporting a padded graph is the same as choosing
that transported final allowance directly. This is literal graph equality.


\lean{Copying/ObservationPadding}{9}{15}{ObservationPadding.padTo_add}

\lean{Copying/ObservationPadding}{16}{32}{ObservationPadding.padTo_trans}

An open program consists of its graph together with focus-free and hole-free
proofs. Pad its graph either by a given number of steps or to a specified
adequate allowance, and transport both certificates using the established
label-property laws. Equality of the underlying graphs implies equality of
these packaged programs, since the other fields are proofs. Consequently
padding an open program to $d+n$ agrees with padding it by $n$ steps.


\lean{Copying/ObservationPadding}{33}{37}{ObservationPadding.pad}

\lean{Copying/ObservationPadding}{38}{43}{ObservationPadding.padTo}

\lean{Copying/ObservationPadding}{44}{46}{ObservationPadding.ext_graph}

\lean{Copying/ObservationPadding}{47}{53}{ObservationPadding.padTo_add}

For a closed program, pad its stored graph and increase its stored depth by
the same number of steps. Focus-freedom and hole-freedom again follow from
padding preservation. Thus every padded program remains an admissible input
to the existing execution initializer.


\lean{Copying/ObservationPadding}{54}{62}{ObservationPadding.pad}

A closing context with local allowance $d$ and host allowance $D$ lifts to
one with allowances $d+n$ and $D+n$. Pad its actual host graph, take the
previously constructed padded site, and retain its sole-hole and binding-free
certificates using their proved padding laws. This construction works for
every original closing context, including holes below arbitrary boxes; no
new restriction on observation contexts is imposed.

Closing a padded program in this padded context is recursively isomorphic to
padding the original closed program. This is exactly the iterated deep-filling
law, now with the closing-context and program certificates packaged around
it. The isomorphism compares graphs at the same resulting host allowance.


\lean{Copying/ObservationPadding}{63}{70}{ObservationPadding.pad}

\lean{Copying/ObservationPadding}{71}{80}{ObservationPadding.pad_close_iso}

\subsection{The execution obligation and observation reflection}
\termdef{copying-execution-padding}

Define execution padding invariance by requiring, for every admissible closed
program $P$, universal-cost machine equivalence between its initial state and
the initial state of $P$ padded once. Universal cost forgets the number of
steps but still observes successful termination, not stuckness. This property
is explicitly passed to every theorem that uses it.

The one-step property implies invariance under any finite number of steps of
padding. Induct on their number. Zero padding gives reflexive refinement in
both directions. At a successor, apply the one-step obligation to the already
padded program and compose each direction with the induction hypothesis. The
original iterative definition of padding makes the final program coincide
with the desired successor padding.


\lean{Copying/ObservationPadding}{81}{84}{ObservationPadding.RespectsPadding}

\lean{Copying/ObservationPadding}{85}{94}{ObservationPadding.initial_pad}

Let $C$ be a closing context and $G$ an open program. First apply iterated
execution padding invariance to the closed program $C[G]$. Next apply the
execution's separate isomorphism obligation to the proved isomorphism between
its padded graph and the graph obtained by closing the padded $G$ in the
padded $C$. Compose the two machine equivalences, reversing their order in
the converse direction. This relates the two actual initialized closed
observations, including their admissibility proofs.


\lean{Copying/ObservationPadding}{95}{106}{ObservationPadding.close_pad}

Suppose the two padded open programs refine at allowance $d+n$. To prove
refinement at $d$, take any original closing context $C$. Transfer the first
closed observation to the padded context, apply the supplied refinement there,
and transfer the second observation back. Transitivity gives the requested
refinement in $C$. Since $C$ was arbitrary, the theorem covers the original
full observation class. Apply this argument in both directions to obtain
reflection of equivalence.

For a specified adequate allowance $N\geq d$, write $N=d+n$ and use the exact
open-program padding equality to reduce to this result. These are reflection
laws from larger allowances. No converse, no automatic equivalence between
all representations, and no copying execution is inferred from them.


\lean{Copying/ObservationPadding}{107}{113}{ObservationPadding.refines_of_pad}

\lean{Copying/ObservationPadding}{114}{118}{ObservationPadding.equiv_of_pad}

\lean{Copying/ObservationPadding}{119}{128}{ObservationPadding.equiv_of_padTo}

Recursive isomorphism of open program graphs also gives observational
equivalence when execution respects isomorphism of closed graphs. In each
closing context, deep replacement congruence lifts the fragment isomorphism
to an isomorphism of the closed fillings. Apply the closed execution obligation
and select its forward and converse refinements for the two contextual
directions. Padding invariance is unnecessary for this equal-allowance law.


\lean{Copying/ObservationPadding}{129}{135}{ObservationPadding.equiv_iso}

\subsection{A cofinal criterion for the unchanged term equivalence}
\termdef{copying-cofinal-observations}

The translated-program wrapper uses the original term translation padded to
the chosen allowance. Padding that wrapper through an intermediate adequate
allowance gives exactly the wrapper at the final allowance. Apply graph-level
padding transitivity and then the open-program equality principle. The source
term, its translation and its interfaces are unchanged.


\lean{Copying/ObservationPadding}{136}{143}{ObservationPadding.translatedProgram_padTo}

Under the two execution representation laws, the original definition
$\mathsf{ObservationalEquiv}(t,u)$ is equivalent to:
\[
 \forall b\;\exists N\geq b,\quad
 \mathsf{depth}(t)\leq N\ \land\ \mathsf{depth}(u)\leq N\ \land\
 \mathsf{Equiv}_{\mathrm{any}}(T_N(t),T_N(u)),
\]
where $T_N$ is the unchanged translated-program wrapper. The witnessing
allowances need not be every natural number; they must be arbitrarily large.

For the forward direction, given $b$, choose the maximum of $b$ and the two
term depths, then instantiate the original all-allowances hypothesis. For the
converse, fix any adequate original allowance $d$. Obtain a witnessing
$N\geq d$ from the cofinal hypothesis. The translated-program transitivity
law identifies the translations at $N$ with the translations at $d$ padded
to $N$. Equivalence reflection therefore supplies the required comparison
at $d$. This argument handles different minimal term depths and preserves
quantification over all original binding-free closing contexts.

This reduction supports later common-allowance proofs. It does not prove the
cofinal hypothesis for Proposition~12.2, or remove the need to establish the
two representation laws for its concrete machine.


\lean{Copying/ObservationPadding}{144}{160}{ObservationPadding.observationalEquiv_iff_cofinal}

\subsection{A negative regression for the missing machine law}

Consider an artificial execution with natural-number states, no transitions,
and final predicate $d=0$. Its initializer returns the closed graph's stored
allowance. Final states are normal because no transitions exist. The execution
is deterministic for the same reason: an alleged pair of outgoing transitions
already includes a contradiction.

It respects recursive graph isomorphism. Such isomorphisms compare graphs at
the same allowance, and both initializers therefore return exactly the same
state. Reflexive machine refinement supplies both directions, regardless of
the graph's labels or carriers. This execution is deliberately not the copying
UAM; it tests which laws the abstract observation interface actually implies.


\lean{Copying/ObservationPaddingExamples}{10}{17}{ObservationPaddingExamples.depthSensitiveExecution}

\lean{Copying/ObservationPaddingExamples}{18}{22}{ObservationPaddingExamples.depthSensitive_deterministic}

\lean{Copying/ObservationPaddingExamples}{23}{27}{ObservationPaddingExamples.depthSensitive_respectsIso}

Take the unchanged closed unit translation at allowance zero, with its
proved focus-free and hole-free certificates. Its initial state is zero,
so the zero-step run is a successful observation. Pad the same closed program
once. Its initial state is now one. Every run in the transition-free machine
has zero length and unchanged state; it therefore cannot reach a final state.
This proves failure of termination for the padded program.

If the execution satisfied padding invariance, its forward refinement for
this unit program would turn the successful zero-step run into a successful
run from the padded initializer, contradicting the preceding result.
Thus padding invariance is independent of the supplied determinism and
isomorphism proofs. An axiom audit of a conditional transfer theorem cannot
substitute for establishing this missing machine property.

Appendix~B.1 retains observational term-context congruence and the concrete
execution obligations. The common-allowance reduction and explicit context
lift are now proved; the copying UAM, the template proofs and the final
proposition remain unfinished.


\lean{Copying/ObservationPaddingExamples}{28}{32}{ObservationPaddingExamples.zeroDepthUnit}

\lean{Copying/ObservationPaddingExamples}{33}{36}{ObservationPaddingExamples.depthSensitive_unit_terminates}

\lean{Copying/ObservationPaddingExamples}{37}{45}{ObservationPaddingExamples.depthSensitive_padded_unit_not_terminates}

\lean{Copying/ObservationPaddingExamples}{46}{51}{ObservationPaddingExamples.depthSensitive_not_respectsPadding}
